%------------------------------------------------------------------------------
\subsection{\file{FastHydrology.jl}: the module and its exports}\label{ref:module}
\emph{Tour PART 0.}

The module file loads the dependencies (\code{Statistics: mean}, \code{ImageFiltering}, \code{SpecialFunctions: erf},
\code{FFTW}, \code{OffsetArrays}, \code{DocStringExtensions}, \code{MAT}, \code{NCDatasets}), includes the 22 other source files and
declares the exported names grouped by file. It defines nothing itself. There is no compile-time configuration: no
preferences, no backend, no float-type constant (contrast Shakti, whose backend is fixed at load time). Exported
names, by file:
\begin{center}\small
\begin{longtable}{>{\raggedright\arraybackslash}p{4.2cm}>{\raggedright\arraybackslash}p{11.2cm}}
\toprule
file & exported names \\
\midrule
\endhead
\file{common/\allowbreak{}grid.jl} & \code{Abstract\allowbreak{}Hydro\allowbreak{}Grid}, \code{ArrayHydroGrid}, \code{fill\_halo!}, \code{alloc\_field}, \code{convolve!}, \code{masked\_mean},
  \code{masked\_max\_abs}, \code{masked\_max\_abs\_diff}, \code{overwrite\_where!}, \code{minus\_gradient\_x!}, \code{minus\_gradient\_y!} \\
\file{common/\allowbreak{}model.jl} and the model files & \code{Abstract\allowbreak{}Hydro\allowbreak{}Model}, \code{Kazmierczak\allowbreak{}Hydro\allowbreak{}Model}, \code{HABHydroModel}, \code{Shakti\allowbreak{}Hydro\allowbreak{}Model},
  \code{Abstract\allowbreak{}Drainage\allowbreak{}Mode}, \code{BothDrainage}, \code{EfficientOnly}, \code{InefficientOnly}, the sliding laws
  (\code{Abstract\allowbreak{}Sliding\allowbreak{}Law}, \code{No\allowbreak{}Friction\allowbreak{}Sliding\allowbreak{}Law}, \code{Prescribed\allowbreak{}Field\allowbreak{}Sliding\allowbreak{}Law}, \code{Weertman\allowbreak{}Sliding\allowbreak{}Law}, \code{Power\allowbreak{}Plastic\allowbreak{}Sliding\allowbreak{}Law},
  \code{Regularized\allowbreak{}Coulomb\allowbreak{}Sliding\allowbreak{}Law}, \code{Regularized\allowbreak{}Coulomb\allowbreak{}Field\allowbreak{}Sliding\allowbreak{}Law}, \code{Shakti\allowbreak{}Regularized\allowbreak{}Coulomb\allowbreak{}Sliding\allowbreak{}Law}), the psi-out
  algorithms (\code{Abstract\allowbreak{}Psi\allowbreak{}Out\allowbreak{}Algorithm}, \code{RecursivePsiOut}, \code{IterativePsiOut}, \code{Topological\allowbreak{}Psi\allowbreak{}Out}, \code{TapedPsiOut}),
  the fill algorithms (\code{Abstract\allowbreak{}Fill\allowbreak{}Algorithm}, \code{JacobiFill}, \code{Lowest\allowbreak{}Neighbour\allowbreak{}Fill}, \code{Priority\allowbreak{}Flood\allowbreak{}Fill}), the routing
  schemes (\code{Abstract\allowbreak{}Routing\allowbreak{}Scheme}, \code{GDSWarner}, \code{Warner}, \code{Quinn}, \code{Tarboton}, \code{Modified\allowbreak{}Tarboton}, \code{GDSTarboton}),
  \code{Abstract\allowbreak{}QConversion}, \code{QFromOutflow}, \code{QFrom\allowbreak{}Face\allowbreak{}Average}, \code{Abstract\allowbreak{}Dissipation\allowbreak{}Discretization}, \code{Cell\allowbreak{}Centred\allowbreak{}Dissipation},
  \code{FaceDissipation}, \code{Abstract\allowbreak{}Friction\allowbreak{}Discretization}, \code{Cell\allowbreak{}Centred\allowbreak{}Friction}, \code{Staggered\allowbreak{}Friction}, the thickness closures
  (\code{Abstract\allowbreak{}Water\allowbreak{}Thickness\allowbreak{}Algorithm}, \code{Areal\allowbreak{}Conduit\allowbreak{}Thickness}, \code{Darcy\allowbreak{}Weisbach\allowbreak{}Thickness}, \code{Laminar\allowbreak{}Thickness}),
  \code{Abstract\allowbreak{}Gradient\allowbreak{}Convention}, \code{LocalGradient}, \code{MeanGradient} \\
\file{models/\allowbreak{}kazmierczak2024/\allowbreak{}sliding\_law.jl}, \file{routing.jl} & \code{calc\_tau\_b}, \code{update\_tau\_b!}, \code{set\_basal\_terms!},
  \code{freeze\_on\_capacity!} \\
\file{common/\allowbreak{}state.jl}, \file{simulation.jl}, \file{run.jl} & \code{Abstract\allowbreak{}Hydro\allowbreak{}State}, \code{HydroState}, \code{Abstract\allowbreak{}Simulation},
  \code{TimeSimulation}, \code{Steady\allowbreak{}State\allowbreak{}Simulation}, \code{run!}, \code{step!}, \code{update\_steady\_state!} \\
\file{models/\allowbreak{}kazmierczak2024/\allowbreak{}water\_flux.jl} & \code{update\_q!}, \code{update\_W!}, \code{update\_phi0!}, \code{potential\_filling!},
  \code{update\_potential\_gradients!}, \code{update\_smoothed\_potential\_gradients!}, \code{accumulate\_psi\_out!}, \code{update\_psi\_out!},
  \code{update\_psi\_out\_iterative!}, \code{update\_psi\_out\_taped!}, \code{route\_psi\_out!}, \code{invalidate\_routing\_tape!} \\
effective pressure files & \code{update\_N!}, \code{update\_Po!}, \code{update\_p\_w!}, \code{N\_responds\_to\_ub}, \code{N\_from\_ub!}, \code{update\_H!},
  \code{update\_S\_inf!}, \code{update\_N\_inf!}, \code{update\_Q!} \\
\file{models/\allowbreak{}kazmierczak2024/\allowbreak{}data\_loaders.jl} & \code{load\_Kazmierczak}, \code{load\_yelmox}, \code{load\_Kazmierczak\_melt} \\
\file{common/\allowbreak{}utilities.jl}, \file{plotting.jl} & \code{compute\_lims}, \code{perYear2perSecond}, \code{perSecond2perYear}, \code{Km2m},
  \code{visualize\_grid}, \code{visualize\_field}, \code{mask\_field} \\
\file{common/\allowbreak{}checkpoint.jl}, \file{output.jl} & \code{save\_checkpoint}, \code{load\_checkpoint!}, \code{Abstract\allowbreak{}Output\allowbreak{}Writer},
  \code{Net\allowbreak{}CDFOutput\allowbreak{}Writer}, \code{write\_output!}, \code{close\_output!} \\
\bottomrule
\end{longtable}
\end{center}
Not exported but used by advanced callers: \code{KazmierczakParams}, \code{KazmierczakWorkspace}, \code{RoutingTape},
\code{DissipationMeltOn/Off}, \code{AbstractPressureDependentSlidingLaw}, \code{update\_psi\_out\_topological!}, \code{resolve\_q!},
\code{prepare\_routing!}, \code{KAZMIERCZAK\_DEFAULT\_L\_W}. The tour checks that every exported name is defined.

%------------------------------------------------------------------------------
\subsection{\file{common/grid.jl}: the grid interface and \code{ArrayHydroGrid}}\label{ref:grid}
\emph{Tour PART 1.}

The interface every model is written against (Section~\ref{arch:types}), its plain-array defaults, and the one grid.
\begin{fnlist}
\item[AbstractHydroGrid] Supertype of grids. A concrete grid must have fields \code{Nx}, \code{Ny}, \code{dx}, \code{dy}.
\item[ArrayHydroGrid\{T\}(Nx, Ny, dx, dy, conv\_cache)] The only grid. Immutable; \code{conv\_cache::Base.RefValue\{Any\}}
  holds the \code{FFTConvCache} of \code{convolve!}, a \code{Ref} so the struct stays immutable while its contents are built on
  first use.
\item[ArrayHydroGrid(Nx, Ny, xlims, ylims; T = Float64)] The constructor: $\Dx = T(x_2-x_1)/N_x$ (cell-centred, limits are
  edges); \code{Nx}, \code{Ny} $>0$ or \code{ArgumentError}.
\item[alloc\_field(grid) / alloc\_field(grid, data)] The only place that knows the field type. For \code{ArrayHydroGrid}:
  \code{zeros(T, Nx, Ny)}, and \code{T.(reshape(data, Nx, Ny))} (a converted copy). The generic methods throw an error naming the
  grid type. Every array of a state and a model is allocated through it.
\item[fill\_halo!(field, grid)] A no-op. Halo-carrying grids would override it; the physics still calls it after writing
  some fields (a v1 remnant, harmless).
\item[convolve!(grid, dest, src, kernel)] Generic method: \code{imfilter!(dest, src, centered(kernel))} (replicate border),
  which allocates new FFT plans and buffers on every call. \code{ArrayHydroGrid} method:
  \code{cached\_fft\_convolve!(grid.conv\_cache, ...)}. The two agree to $10^{-12}$ (tour).
\item[minus\_gradient\_x\_kernel!(dest, field, Nx, Ny, dx), minus\_gradient\_y\_kernel!] The loops of
  Section~\ref{disc:gradients}: centred inside, full one-sided on the edges, 0 for a single column/row. Shared by every
  grid so the edge rule cannot drift.
\item[minus\_gradient\_x!(grid, dest, field), minus\_gradient\_y!] Call the kernels with the grid's sizes.
\item[masked\_mean(grid, field, mask)] Mean over \code{mask == 1}, in one branchless pass (\code{ifelse} with \code{@simd}); the
  source records this beat collecting the grounded indices (\code{findall}) at every grounded fraction from 1\,\% to 95\,\%.
  Returns $0/0$ = NaN for an empty mask.
\item[masked\_max\_abs(grid, field, mask), masked\_max\_abs\_diff(grid, a, b, mask)] $\max|f|$ and $\max|a-b|$ over \code{mask == 1},
  same style; used for the relative convergence tests of the Picard loops.
\item[overwrite\_where!(grid, dest, cond, predicate, src; scale = true)] \code{dest[i] = scale*src} (scalar \code{src}) or
  \code{scale*src[i]} (field) where \code{predicate(cond[i])}, in one branchless pass instead of materialising
  \code{predicate.(cond)}. Used to resolve the degenerate cases of $\Sinf$, $H_{\rm soft}$, $\Ninf$ and by HAB's $p_w$.
\end{fnlist}

%------------------------------------------------------------------------------
\subsection{\file{common/fft\_convolution.jl}: a plan-reusing FFT convolution}\label{ref:fft}
\emph{Tour PART 2.}

The smoothing of Section~\ref{disc:smoothing} convolves two fields (or one potential) with the same kernel on every
solve, and a coupled run solves thousands of times on the same grid. ImageFiltering's FFT path builds FFTW plans and
padded complex buffers on every call (the source measured ~20\,MB per call on a real grid) and exposes no way to keep
them. This file keeps them. Only point-symmetric kernels are supported (then convolution equals correlation); the
kernel array may be rectangular.
\begin{fnlist}
\item[FFTConvCache] Mutable struct: kernel radii \code{frb\_x}, \code{frb\_y}; the padded size \code{full\_size}; buffers
  \code{padded\_src}, \code{padded\_kernel}, \code{result\_full} (\code{Matrix\{Float64\}}), \code{Fsrc}, \code{Fkern}, \code{Fresult}
  (\code{Matrix\{ComplexF64\}}); the FFTW plans \code{plan}, \code{inv\_plan} (typed \code{Any}, to avoid naming FFTW's internal plan
  types); and \code{kernel}, a copy of the kernel \code{Fkern} was computed from.
\item[fft\_size(n)] \code{nextprod((2, 3, 5, 7), n)}: the smallest 7-smooth size $\ge n$. FFTW is slow for sizes with a large prime
  factor (Thwaites 2\,km padded to $389\times332$, 389 prime). Padding further does not change the cropped result: an
  output cell only reads inputs within the kernel radius, all inside the first $N_x + 2r_x$ rows, so nothing wraps around.
\item[FFTConvCache(Nx, Ny, frb\_x, frb\_y)] Allocates the buffers and builds the plans (\code{plan\_rfft}, and \code{plan\_irfft}
  with the original row count passed explicitly, since a real FFT's output cannot tell an even length from an odd one).
  \code{kernel} starts as NaN, which never compares equal, so the first call always transforms the kernel.
\item[get\_fft\_conv\_cache!(cache\_ref, Nx, Ny, frb\_x, frb\_y)] Returns the cached object, rebuilding it on first use or when
  the radii or padded size changed. Asserts the concrete type on return so callers stay type-stable although the
  \code{Ref} is \code{Any}.
\item[fill\_padded\_src!(padded, src, frb\_x, frb\_y)] Data in the middle; edge rows and columns replicated outward over the
  border; corners from the corner pixels (ImageFiltering's default ``replicate'' border).
\item[fill\_padded\_kernel!(padded, kernel2d, frb\_x, frb\_y, full\_size)] Places kernel offset $(d_i, d_j)$ at
  \code{[mod(di, Nfx)+1, mod(dj, Nfy)+1]}: centre at \code{[1, 1]}, negative offsets wrapped to the far end, so a circular FFT
  convolution lines up with the padding.
\item[cached\_fft\_convolve!(cache\_ref, dest, src, kernel)] Pads \code{src}, forward FFT, re-transforms the kernel only if it
  changed (the same kernel serves both gradient components and, in a coupled run, every step while the mean thickness is
  unchanged), multiplies, inverse FFT, crops into \code{dest}. A warm call allocates (almost) nothing (tour). \code{dest} and
  \code{src} are used as given: a halo-carrying grid would have to pass interior views.
\end{fnlist}

%------------------------------------------------------------------------------
\subsection{\file{common/model.jl}: \code{AbstractHydroModel}}\label{ref:model}
\emph{Tour PART 3.}
\begin{fnlist}
\item[AbstractHydroModel] Supertype of the three models.
\item[N\_responds\_to\_ub(model)] Default \code{false}. Whether $N$ depends on the basal sliding speed, so that an ice model whose
  velocity solve depends on $N$ should re-evaluate $N$ inside that solve (\code{N\_from\_ub!}) rather than once per step. The
  docstring gives the reason: a steady hydrology evaluated once per step from the previous velocity makes $N$ and $u_b$
  alternate between two states. HAB (no $u_b$ dependence) and Shakti ($N$ is prognostic) keep \code{false}; K24 overrides it
  (\ref{ref:k24ep}).
\end{fnlist}

%------------------------------------------------------------------------------
\subsection{\file{common/state.jl}: \code{HydroState}}\label{ref:state}
\emph{Tour PART 4.}
\begin{fnlist}
\item[AbstractHydroState] Supertype of states.
\item[HydroState\{A\}(mask, h, b, N, W)] Inputs \code{mask} (1 = grounded), \code{h} [m], \code{b} [m]; outputs \code{N} [Pa], \code{W} [m].
  Immutable: models write into the arrays.
\item[HydroState(grid, mask\_in, h\_in, b\_in)] Checks each input is \code{(Nx, Ny)} (an \code{ArgumentError} names the field), copies
  them through \code{alloc\_field} (so the state never aliases the caller's arrays, and an integer mask becomes the grid's float
  type), and allocates \code{N}, \code{W} as zeros.
\end{fnlist}

%------------------------------------------------------------------------------
\subsection{\file{common/simulation.jl}: the simulation types}\label{ref:simulation}
\emph{Tour PART 5.}
\begin{fnlist}
\item[AbstractSimulation] Supertype; each subtype has a \code{run!}.
\item[SteadyStateSimulation\{M, G, S\}(model, grid, state)] For K24 and HAB. No time, no loop.
\item[TimeSimulation\{M\}(model)] Wraps a model that owns its time evolution (today only \code{ShaktiHydroModel}); its
  \code{run!}/\code{step!} methods come from the extension.
\end{fnlist}

%------------------------------------------------------------------------------
\subsection{\file{common/effective\_pressure.jl}: the overburden}\label{ref:commonep}
\emph{Tour PART 6.}
\begin{fnlist}
\item[update\_Po!(model, grid, state)] \code{model.Po .= model.rho\_i * model.g * state.h}, for any model with those fields. HAB calls
  it; K24 computes $p_o$ inside its fused kernel, with the same expression.
\end{fnlist}

%------------------------------------------------------------------------------
\subsection{\file{common/run.jl}: \code{run!} and \code{step!}}\label{ref:run}
\emph{Tour PART 7.}
\begin{fnlist}
\item[run!(sim::SteadyStateSimulation)] \code{update\_steady\_state!(sim.model, sim.grid, sim.state)}.
\item[step!] Declared with no methods (\code{function step! end}). The Shakti extension adds the only method. It exists so a coupled
  driver owning its own time loop can advance a time-evolving model one step at a time.
\end{fnlist}

%------------------------------------------------------------------------------
\subsection{\file{common/checkpoint.jl}: save and restore}\label{ref:checkpoint}
\emph{Tour PART 8.}

K24 and HAB have no time-integrated state: every \code{run!} recomputes $N$ and $W$ from scratch. Checkpoints are for the
coupled driver: on resuming a killed job it restores the hydrology side (the geometry the checkpointed step used, and
$N$, $W$) together with its own ice state. The format is NetCDF, since NCDatasets is already a dependency (Shakti uses
JLD2).
\begin{fnlist}
\item[\_output\_array(field)] \code{Array(field)}: a plain copy ready to write. The hook a non-Array grid would extend. Shared with
  \code{write\_output!}.
\item[\_copy\_into!(field, data)] \code{field .= data}. Inverse hook.
\item[save\_checkpoint(path, state, step, time)] Writes every field of any \code{AbstractHydroState} (by \code{fieldnames}) as an
  \code{(xc, yc)} \code{Float64} variable, with \code{step}/\code{time} attributes, to \code{path.tmp}, then \code{mv}s it over \code{path}. An atomic
  rename on the same file system: a crash during the write leaves the previous checkpoint intact. A state with no fields
  is an error.
\item[load\_checkpoint!(state, path)] Restores every field in place and returns \code{(step::Int, time::Float64)}.
\end{fnlist}

%------------------------------------------------------------------------------
\subsection{\file{common/output.jl}: periodic NetCDF output}\label{ref:output}
\emph{Tour PART 9.}

The motivating incident, recorded in the file header: a ~4000-step coupled Yelmo--FastHydrology run (Kryonomos MISMIP) that
wrote its NetCDF result only at the end lost everything at the wall-time limit. The writer appends one slice per call
and syncs it to disk. Shakti has its own output machinery, so this exists for the steady models.
\begin{fnlist}
\item[AbstractOutputWriter] Supertype; dispatch on it picks the format.
\item[NetCDFOutputWriter\{G\}(path, grid, field\_names, ds)] \code{ds} is the open \code{NCDataset} (typed \code{Any}).
\item[NetCDFOutputWriter(path, grid, field\_names; resume\_step = nothing, overwrite = true)] Without \code{resume\_step}: creates the
  file (dimensions \code{xc}, \code{yc}, unlimited \code{time}; variables \code{step}, \code{time} and one \code{(xc, yc, time)} variable per field),
  deleting an existing file if \code{overwrite} or erroring otherwise. With \code{resume\_step = k}: the file must exist; slices with
  \code{step > k} are dropped (\code{\_trim\_output\_after!}), it is opened for appending, and every field must exist in it.
\item[\_trim\_output\_after!(path, field\_names, resume\_step)] A NetCDF unlimited dimension can grow but not shrink, so trimming copies
  the kept slices to \code{path.trim.tmp} and renames. Writes with explicit index ranges (\code{v[1:n] = ...}): the source records
  that \code{v[:] = ...} on a still-empty unlimited dimension is a silent no-op in NCDatasets.
\item[write\_output!(writer, step, time, fields::NamedTuple)] Appends one slice: \code{fields[name]} for each of the writer's
  names (extra entries ignored, a missing one an error, raised after \code{step}/\code{time} are already appended:
  Section~\ref{pit:open}), then \code{NCDatasets.sync}.
\item[close\_output!(writer)] Closes the dataset.
\end{fnlist}

%------------------------------------------------------------------------------
\subsection{\file{common/utilities.jl}: units and limits}\label{ref:utilities}
\emph{Tour PART 10.}
\begin{fnlist}
\item[SECONDS\_PER\_YEAR, METERS\_PER\_KILOMETER] $365.25\times86400$ (Julian year) and 1000.
\item[perYear2perSecond(f), perSecond2perYear(f), Km2m(f)] Divide or multiply in the argument's float type.
\item[compute\_lims(xc, yc)] Edge limits from cell-centre coordinates: \code{extrema} padded by half the spacing of the first two
  entries.
\end{fnlist}

%------------------------------------------------------------------------------
\subsection{\file{common/plotting.jl}: plotting entry points}\label{ref:plotting}
\emph{Tour PART 11.}
\begin{fnlist}
\item[visualize\_field, visualize\_grid] Declared without methods; the Makie extension implements them (\ref{ref:makieext}).
\item[mask\_field(field, mask, value)] A copy of \code{field} with every cell where \code{mask != 1} set to \code{value} (NaN blanks the
  ocean in a heatmap). Needs no Makie, so it lives here.
\end{fnlist}
